\documentclass[11pt,a4paper]{article}
\usepackage[margin=2.5cm]{geometry}
\usepackage{amsmath}
\usepackage{booktabs}
\usepackage{xcolor}
\definecolor{accent}{RGB}{36,87,166}
\title{Single-cell conductance model\\\large Equations for the Brian2 tutorial}
\author{}
\date{}
\begin{document}
\maketitle

\section*{Overview}
The tutorial uses one Hodgkin--Huxley neuron with a fast sodium conductance,
a delayed-rectifier potassium conductance, and a passive leak conductance. The
leak reversal potential is deliberately depolarized to $-30\,\mathrm{mV}$.
Thus, at the initial voltage ($-65\,\mathrm{mV}$), the leak current is inward
and the neuron fires without an additional injected current.

\section*{Membrane equation}
Using the convention that positive current depolarizes the membrane,
\begin{align}
C_m\frac{dV}{dt} &= I_{\mathrm{Na}}+I_{\mathrm{K}}+I_{\mathrm{L}},\\
I_{\mathrm{Na}} &= \bar g_{\mathrm{Na}}m^3h(E_{\mathrm{Na}}-V),\\
I_{\mathrm{K}} &= \bar g_{\mathrm{K}}n^4(E_{\mathrm{K}}-V),\\
I_{\mathrm{L}} &= g_{\mathrm{L}}(E_{\mathrm{L}}-V).
\end{align}
The variables $m$ and $h$ describe sodium activation and inactivation;
$n$ describes potassium activation.

\section*{Gating variables}
For each gate $x\in\{m,h,n\}$,
\begin{equation}
\frac{dx}{dt}=\alpha_x(V)(1-x)-\beta_x(V)x.
\end{equation}
With $V$ measured in mV and all rates in ms$^{-1}$,
\begin{align}
\alpha_m(V)&=\frac{(V+40)/10}{1-\exp[-(V+40)/10]}, &
\beta_m(V)&=4\exp[-(V+65)/18],\\
\alpha_h(V)&=0.07\exp[-(V+65)/20], &
\beta_h(V)&=\frac{1}{1+\exp[-(V+35)/10]},\\
\alpha_n(V)&=0.1\frac{(V+55)/10}{1-\exp[-(V+55)/10]}, &
\beta_n(V)&=0.125\exp[-(V+65)/80].
\end{align}
Brian2 evaluates the removable singularities in $\alpha_m$ and $\alpha_n$
with its numerically stable \texttt{exprel} function.

\section*{Parameters}
\begin{center}
\begin{tabular}{lll}
\toprule
Parameter & Value & Meaning\\
\midrule
$C_m$ & $1\,\mu\mathrm{F}/\mathrm{cm}^2$ & membrane capacitance density\\
$\bar g_{\mathrm{Na}}$ & $120\,\mathrm{mS}/\mathrm{cm}^2$ & maximal sodium conductance\\
$\bar g_{\mathrm{K}}$ & $36\,\mathrm{mS}/\mathrm{cm}^2$ & maximal potassium conductance\\
$g_{\mathrm{L}}$ & $0.3\,\mathrm{mS}/\mathrm{cm}^2$ & leak conductance\\
$E_{\mathrm{Na}}$ & $50\,\mathrm{mV}$ & sodium reversal potential\\
$E_{\mathrm{K}}$ & $-77\,\mathrm{mV}$ & potassium reversal potential\\
$E_{\mathrm{L}}$ & $-30\,\mathrm{mV}$ & depolarizing leak reversal potential\\
\bottomrule
\end{tabular}
\end{center}

\section*{Initial conditions and spike detection}
The membrane starts at $V(0)=-65\,\mathrm{mV}$. Each gate starts at its
steady-state value,
\begin{equation}
x(0)=x_\infty(V(0))=\frac{\alpha_x(V(0))}
{\alpha_x(V(0))+\beta_x(V(0))}.
\end{equation}
A spike is counted when $V$ crosses the operational sodium-activation threshold
of $-55\,\mathrm{mV}$ upward. The Hodgkin--Huxley sodium gate is continuous, so
this is a pedagogical measurement threshold rather than a hard channel switch.
Action potentials are generated by the conductance dynamics and voltage is
never reset.

\end{document}
